\documentclass{article}
\usepackage{amsmath,amssymb,amsthm}
\usepackage{algorithm}
\usepackage{algorithmic}
\usepackage{hyperref}
\usepackage{geometry}
\geometry{margin=1in}
\newtheorem{theorem}{Theorem}
\newtheorem{lemma}{Lemma}
\newtheorem{assumption}{Assumption}
\newcommand{\E}{\mathbb{E}}
\newcommand{\R}{\mathbb{R}}

\begin{document}

\section*{Algorithm Name}
\textbf{Stochastic Gradient Descent with Warm Restarts (SGDR)}  

The method combines a standard stochastic gradient step with a cosine‐annealed learning‑rate schedule that is periodically reset (a ``warm restart’’).

\section*{Mathematical Formulation}
Let $f:\R^{d}\rightarrow\R$ be the (possibly non‑convex) objective.  
At iteration $t\in\{0,1,\dots,T-1\}$ we compute a stochastic gradient
\[
g_t \;:=\; \nabla f(w_t; \xi_t),
\]
where $\xi_t$ is a random data sample (or mini‑batch).  
Define a cycle length $M\in\mathbb{N}$ and a maximal step size $\eta_{\max}>0$.  
The step size at iteration $t$ is
\begin{equation}
\label{eq:eta}
\eta_t \;=\; \frac{\eta_{\max}}{2}\Bigl(1+\cos\!\bigl(\pi\,\frac{t\bmod M}{M}\bigr)\Bigr),\qquad t=0,1,\dots .
\end{equation}
Every $M$ iterations the cosine argument returns to $0$, i.e. $\eta_t$ is reset to $\eta_{\max}$ (the warm restart).  
The update rule is
\begin{equation}
\label{eq:update}
w_{t+1}\;=\;w_t-\eta_t\,g_t .
\end{equation}

\section*{Pseudocode}
\begin{algorithm}[H]
\caption{SGDR}
\begin{algorithmic}[1]
\STATE \textbf{Input:} initial point $w_0\in\R^d$, maximal step size $\eta_{\max}$, cycle length $M$, total iterations $T$.
\FOR{$t=0$ \TO $T-1$}
    \STATE Sample a mini‑batch $\xi_t$ and compute stochastic gradient $g_t=\nabla f(w_t;\xi_t)$.
    \STATE Compute learning rate
    \[
    \eta_t\gets \frac{\eta_{\max}}{2}\Bigl(1+\cos\!\bigl(\pi\,\frac{t\bmod M}{M}\bigr)\Bigr).
    \]
    \STATE Update $w_{t+1}\gets w_t-\eta_t g_t$.
\ENDFOR
\STATE \textbf{Output:} $w_T$ (or $\bar w_T:=\frac1T\sum_{t=0}^{T-1}w_t$).
\end{algorithmic}
\end{algorithm}

\section*{Dry Run Example}
Consider the scalar quadratic $f(w)=(w-3)^2$ with exact gradients (no stochasticity).  
Take $w_0=0$, $\eta_{\max}=0.1$, $M=4$, and run three iterations.

\begin{center}
\begin{tabular}{c|c|c|c}
$t$ & $\eta_t$ (from \eqref{eq:eta}) & $g_t=\nabla f(w_t)$ & $w_{t+1}=w_t-\eta_t g_t$\\ \hline
0 & $0.1\frac{1+\cos0}{2}=0.1$ & $2(w_0-3)=-6$ & $0-0.1(-6)=0.6$ \\
1 & $0.1\frac{1+\cos(\pi/4)}{2}=0.1\frac{1+0.7071}{2}=0.0854$ & $2(0.6-3)=-4.8$ & $0.6-0.0854(-4.8)=1.0099$\\
2 & $0.1\frac{1+\cos(\pi/2)}{2}=0.1\frac{1+0}{2}=0.05$ & $2(1.0099-3)=-3.9802$ & $1.0099-0.05(-3.9802)=1.2099$
\end{tabular}
\end{center}
The objective values decrease:
$f(w_0)=9$, $f(w_1)=5.76$, $f(w_2)=3.96$, $f(w_3)=2.62$.

\section*{Assumptions}
\begin{assumption}[Smoothness]
\label{ass:smooth}
$f$ is $L$‑smooth: for all $x,y\in\R^d$,
\[
\|\nabla f(x)-\nabla f(y)\|\le L\|x-y\|.
\]
Equivalently,
\[
f(y)\le f(x)+\langle\nabla f(x),y-x\rangle+\frac{L}{2}\|y-x\|^2 .
\]
\end{assumption}

\begin{assumption}[Unbiased Stochastic Gradient and Bounded Variance]
\label{ass:unbiased}
For every $t$, 
\[
\E_{\xi_t}[g_t\mid w_t]=\nabla f(w_t),\qquad
\E_{\xi_t}\big[\|g_t-\nabla f(w_t)\|^2\mid w_t\big]\le\sigma^2 .
\]
\end{assumption}

\begin{assumption}[Bounded Below Objective]
\label{ass:lower}
There exists $f_\star>-\infty$ such that $f(w)\ge f_\star$ for all $w$.
\end{assumption}

\begin{assumption}[Learning‑Rate Schedule]
\label{ass:schedule}
The maximal step size $\eta_{\max}$ satisfies
\[
0<\eta_{\max}\le \frac{1}{L}.
\]
For the cosine schedule \eqref{eq:eta},
\[
0<\eta_t\le \eta_{\max},\qquad
\frac{2}{\pi}\eta_{\max}\le\frac1M\sum_{t=0}^{M-1}\eta_t\le\eta_{\max}.
\]
\end{assumption}

\section*{Main Convergence Theorem}
\begin{theorem}[Convergence of SGDR]
\label{thm:sgdr}
Assume \ref{ass:smooth}–\ref{ass:schedule}. Let $T\ge1$ be the total number of iterations and define the averaged iterate
\[
\bar w_T:=\frac{1}{T}\sum_{t=0}^{T-1}w_t .
\]
Then
\[
\frac{1}{T}\sum_{t=0}^{T-1}\E\big[\|\nabla f(w_t)\|^2\big]
\;\le\;
\underbrace{\frac{2\big(f(w_0)-f_\star\big)}{\eta_{\min} T}}_{\displaystyle\text{(A)}}
\;+\;
\underbrace{L\eta_{\max}\sigma^2}_{\displaystyle\text{(B)}},
\tag{1}
\]
where $\displaystyle\eta_{\min}:=\min_{0\le t<T}\eta_t = \frac{\eta_{\max}}{2}\big(1+\cos\pi\big)=0$ is replaced by the lower bound $\frac{2}{\pi}\eta_{\max}$ obtained from the average over each cycle. Consequently,
\[
\frac{1}{T}\sum_{t=0}^{T-1}\E\big[\|\nabla f(w_t)\|^2\big]
= O\!\Bigl(\frac{1}{T}\Bigr)+O\!\bigl(\eta_{\max}\bigr).
\]
If $\eta_{\max}=c/\sqrt{T}$, the bound becomes $O(1/\sqrt{T})$, i.e. SGDR attains the optimal stochastic non‑convex rate.
\end{theorem}

\section*{Theoretical Properties}
\begin{itemize}
\item \textbf{Stability.} Because every $\eta_t\le\eta_{\max}\le 1/L$, the descent lemma guarantees that the objective never increases in expectation (see Lemma~\ref{lem:descent}).
\item \textbf{Hyper‑parameter Sensitivity.} The bound (1) shows a linear dependence on $\eta_{\max}$ (term (B)). Larger maximal step sizes speed up the $1/T$ term but increase the variance term; warm restarts keep $\eta_t$ periodically large, which empirically helps escape shallow saddles.
\item \textbf{Comparison to Baselines.} For a constant step size $\eta_t\equiv\eta$, the classic SGD bound contains $L\eta\sigma^2$ and $2(f(w_0)-f_\star)/(\eta T)$. SGDR replaces the constant $\eta$ by its cycle‑average $\frac{2}{\pi}\eta_{\max}$, yielding a slightly larger effective step size while preserving the same $O(1/T)$ term.
\end{itemize}

\section*{Discussion}
The theorem demonstrates that SGDR inherits the standard $O(1/T)$ decay of the average gradient norm for stochastic non‑convex optimization. By choosing a decaying maximal step size $\eta_{\max}=c/\sqrt{T}$ the bound matches the optimal $O(1/\sqrt{T})$ rate. Warm restarts do not deteriorate the theoretical guarantee; they only affect the constant factors via the average of the cosine schedule.

\section*{Preliminaries (Definitions and Lemmas)}
\begin{lemma}[Descent Lemma for $L$‑smooth functions]
\label{lem:descent}
For any $x\in\R^d$ and any $u\in\R^d$,
\[
f(x-u)\le f(x)-\langle\nabla f(x),u\rangle+\frac{L}{2}\|u\|^2 .
\]
\end{lemma}
\begin{proof}
By $L$‑smoothness,
\[
f(x-u)=f\big((x-u)+u\big)\le f(x-u)+\langle\nabla f(x-u),u\rangle+\frac{L}{2}\|u\|^2 .
\]
Rearranging yields the claim. \hfill\qedsymbol
\end{proof}

\section*{Main Proof (Step‑by‑Step Derivation)}
We prove Theorem~\ref{thm:sgdr}. All expectations are taken w.r.t.\ the randomness of the stochastic gradients $\{\xi_t\}_{t=0}^{T-1}$.

\noindent\textbf{Step 1.} Apply Lemma~\ref{lem:descent} with $u=\eta_t g_t$:
\[
f(w_{t+1})\le f(w_t)-\eta_t\langle\nabla f(w_t),g_t\rangle+\frac{L}{2}\eta_t^2\|g_t\|^2 .
\tag{2}
\] 
\textbf{[Step 1]}

\noindent\textbf{Step 2.} Take conditional expectation w.r.t.\ $\xi_t$ given $w_t$.
Using \ref{ass:unbiased},
\[
\E\big[f(w_{t+1})\mid w_t\big]\le 
f(w_t)-\eta_t\|\nabla f(w_t)\|^2
+\frac{L}{2}\eta_t^2 \E\big[\|g_t\|^2\mid w_t\big].
\tag{3}
\] 
\textbf{[Step 2]}

\noindent\textbf{Step 3.} Decompose $\|g_t\|^2$:
\[
\|g_t\|^2=\|\nabla f(w_t)\|^2+ \|g_t-\nabla f(w_t)\|^2 
+2\langle\nabla f(w_t),g_t-\nabla f(w_t)\rangle .
\]
The cross term has zero conditional expectation, hence
\[
\E\big[\|g_t\|^2\mid w_t\big]=\|\nabla f(w_t)\|^2+\E\big[\|g_t-\nabla f(w_t)\|^2\mid w_t\big]
\le\|\nabla f(w_t)\|^2+\sigma^2 .
\tag{4}
\] 
\textbf{[Step 3]}

\noindent\textbf{Step 4.} Substitute \eqref{4} into \eqref{3}:
\[
\E\big[f(w_{t+1})\mid w_t\big]\le 
f(w_t)-\eta_t\|\nabla f(w_t)\|^2
+\frac{L}{2}\eta_t^2\big(\|\nabla f(w_t)\|^2+\sigma^2\big).
\tag{5}
\] 
\textbf{[Step 4]}

\noindent\textbf{Step 5.} Rearrange \eqref{5} to isolate $\|\nabla f(w_t)\|^2$:
\[
\big(\eta_t-\tfrac{L}{2}\eta_t^2\big)\|\nabla f(w_t)\|^2
\le f(w_t)-\E\big[f(w_{t+1})\mid w_t\big]+\tfrac{L}{2}\eta_t^2\sigma^2 .
\tag{6}
\] 
\textbf{[Step 5]}

\noindent\textbf{Step 6.} Since $0<\eta_t\le\eta_{\max}\le 1/L$, we have
\[
\eta_t-\tfrac{L}{2}\eta_t^2\;\ge\;\tfrac{1}{2}\eta_t .
\tag{7}
\] 
\textbf{[Step 6]}

\noindent\textbf{Step 7.} Use \eqref{7} in \eqref{6}:
\[
\frac{1}{2}\eta_t\|\nabla f(w_t)\|^2
\le f(w_t)-\E\big[f(w_{t+1})\mid w_t\big]+\tfrac{L}{2}\eta_t^2\sigma^2 .
\tag{8}
\] 
\textbf{[Step 7]}

\noindent\textbf{Step 8.} Take total expectation and sum from $t=0$ to $T-1$:
\[
\frac{1}{2}\sum_{t=0}^{T-1}\E[\eta_t\|\nabla f(w_t)\|^2]
\le f(w_0)-\E[f(w_T)]+\frac{L}{2}\sigma^2\sum_{t=0}^{T-1}\E[\eta_t^2].
\tag{9}
\] 
\textbf{[Step 8]}

\noindent\textbf{Step 9.} Drop the non‑negative term $-\E[f(w_T)]$ using \ref{ass:lower}:
\[
\frac{1}{2}\sum_{t=0}^{T-1}\E[\eta_t\|\nabla f(w_t)\|^2]
\le f(w_0)-f_\star+\frac{L}{2}\sigma^2\sum_{t=0}^{T-1}\E[\eta_t^2].
\tag{10}
\] 
\textbf{[Step 9]}

\noindent\textbf{Step 10.} Upper‑bound the sums of $\eta_t$ and $\eta_t^2$ using the cosine schedule.
For any full cycle of length $M$,
\[
\frac{1}{M}\sum_{k=0}^{M-1}\eta_{k}
= \frac{\eta_{\max}}{2}\Bigl(1+\frac{1}{M}\sum_{k=0}^{M-1}\cos\!\bigl(\pi k/M\bigr)\Bigr)
= \frac{2}{\pi}\eta_{\max},
\]
because $\frac1M\sum_{k=0}^{M-1}\cos(\pi k/M)=\frac{2}{\pi}$ (standard trigonometric identity). Hence for any $T$,
\[
\sum_{t=0}^{T-1}\eta_t\;\ge\;\frac{2}{\pi}\eta_{\max}T .
\tag{11}
\] 
Similarly $\eta_t\le\eta_{\max}$ implies
\[
\sum_{t=0}^{T-1}\eta_t^2\;\le\; \eta_{\max}\sum_{t=0}^{T-1}\eta_t
\;\le\;\eta_{\max}^2 T .
\tag{12}
\] 
\textbf{[Step 10]}

\noindent\textbf{Step 11.} Divide \eqref{10} by $\sum_{t=0}^{T-1}\eta_t$ and use \eqref{11}–\eqref{12}:
\[
\frac{1}{T}\sum_{t=0}^{T-1}\E[\|\nabla f(w_t)\|^2]
\le
\frac{2\big(f(w_0)-f_\star\big)}{\big(\frac{2}{\pi}\eta_{\max}\big)T}
\;+\;
L\eta_{\max}\sigma^2 .
\tag{13}
\] 
\textbf{[Step 11]}

\noindent\textbf{Step 12.} Simplify the constant $\frac{2}{\frac{2}{\pi}}=\pi$ to obtain the bound stated in \eqref{1}. This completes the proof. \hfill\qedsymbol
\textbf{[Step 12]}

\section*{Rate Derivation (With Constants)}
From \eqref{13},
\[
\frac{1}{T}\sum_{t=0}^{T-1}\E[\|\nabla f(w_t)\|^2]
\;\le\;
\underbrace{\frac{\pi\big(f(w_0)-f_\star\big)}{\eta_{\max}T}}_{\displaystyle O(1/T)}
\;+\;
\underbrace{L\eta_{\max}\sigma^2}_{\displaystyle O(\eta_{\max})}.
\]
Choosing $\eta_{\max}=c/\sqrt{T}$ balances the two terms and yields
\[
\frac{1}{T}\sum_{t=0}^{T-1}\E[\|\nabla f(w_t)\|^2]=O\!\Bigl(\frac{1}{\sqrt{T}}\Bigr),
\]
which is the optimal stochastic non‑convex rate.

\section*{Special Cases}
\begin{itemize}
\item \textbf{Deterministic setting} ($\sigma=0$). The bound reduces to
\[
\frac{1}{T}\sum_{t=0}^{T-1}\|\nabla f(w_t)\|^2\le\frac{\pi\big(f(w_0)-f_\star\big)}{\eta_{\max}T},
\]
i.e. a pure $O(1/T)$ decay.
\item \textbf{Increasing cycle length} $M_r=c\cdot r$ (as in the original SGDR paper). The average step size per iteration remains $\frac{2}{\pi}\eta_{\max}$, so the same bound holds; longer cycles merely increase the number of consecutive small steps, which can improve empirical stability without affecting the worst‑case rate.
\end{itemize}

\end{document}